\subsection{Preliminaries on absolute values}

Let K be a field (commutative or not). Recall (General Topology, Chapter IX, §3, no. 2, Definition 2) that an absolute value on K is any mapping f from K to \(R_{+}\) satisfying the following axioms:

\((\mathrm{VA}_{\mathbf{I}})\) The relation \(f(x) = 0\) is equivalent to \(x = 0\) .

\((\mathrm{VA}_{\mathrm{II}}) f(xy) = f(x)f(y)\) for all \(x, y\) in \(\mathbf{K}\) .

\((\mathrm{VA}_{\mathrm{III}}) f(x + y) \leqslant f(x) + f(y)\) for all \(x, y\) in \(\mathbf{K}\) .

It follows from \((\mathrm{VA}_{\mathbf{I}})\) and \((\mathrm{VA}_{\mathbf{II}})\) that \(f(1) = 1, f(-1) = 1\) and

\[
f (x ^ {- 1}) = \frac {1}{f (x)}
\]

for \(x \neq 0\) .

For a mapping \(f\) from \(\mathbf{K}\) to \(\mathbf{R}_{+}\) and a real number \(\mathbf{A} > 0\) , let \((\mathrm{U}_{\mathbf{A}})\) denote the relation

\[
f (x + y) \leqslant \mathrm{A}. \sup (f (x), f (y)) \quad \text {   for   all   } x, y \text {   in   } \mathrm{K}.
\]

We shall denote by \(\mathcal{V}(\mathbf{K})\) the set of mappings \(f\) from \(\mathbf{K}\) to \(\mathbf{R}_{+}\) satisfying \((\mathrm{VA}_{\mathbf{I}})\) and \((\mathrm{VA}_{\mathbf{II}})\) and for which there exists an \(\mathbf{A} > 0\) (depending on \(f\) ) such that \((\mathbf{U}_{\mathbf{A}})\) holds.

Note that if \(f \in \mathcal{V}(\mathbf{K})\) , then, putting \(x = 1, y = 0\) in \((\mathrm{U}_{\mathbf{A}})\) ,

\[
1 = f (1) \leqslant A. \sup (f (1), f (0)) = A.
\]

PROPOSITION 1. For a mapping \(f\) from \(\mathbf{K}\) to \(\mathbf{R}_+\) satisfying \((\mathrm{VA}_{\mathbf{I}})\) and \((\mathrm{VA}_{\mathbf{II}})\) to belong to \(\mathcal{V}(\mathbf{K})\) , it is necessary and sufficient that \(f(1 + x)\) be bounded in the set of \(x \in \mathbf{K}\) such that \(f(x) \leqslant 1\) .

If \(f\) satisfies \((\mathrm{U}_{\mathbf{A}})\) , then \(f(1 + x) \leqslant \mathbf{A}\) if \(f(x) \leqslant 1\) . Conversely, suppose that \(f(x + 1) \leqslant \mathbf{A}\) for the \(x \in \mathbf{K}\) such that \(f(x) \leqslant 1\) (which implies that \(\mathbf{A} \geqslant f(1) = 1\) ); then, if \(x = 0\) or \(y = 0\) , condition \((\mathrm{U}_{\mathbf{A}})\) is fulfilled; if on the other hand \(x \neq 0\) and \(y \neq 0\) , we may assume for example that \(f(y) \leqslant f(x)\) , hence, by \((\mathrm{VA}_{\mathrm{II}}), f(yx^{-1}) \leqslant 1\) and therefore \(f(1 + yx^{-1}) \leqslant \mathbf{A}\) , which gives, by virtue of \((\mathrm{VA}_{\mathrm{II}}), f(x + y)f(x)^{-1} \leqslant \mathbf{A}\) ; whence

\[
f (x + y) \leqslant A f (x) \leqslant A. \sup (f (x), f (y)).
\]

If \(f\) is an absolute value on \(\mathbf{K}\) , then \(f(n,1) \leqslant n\) by induction on the integer \(n > 0\) starting from \((\mathrm{VA}_{\mathrm{III}})\) ; conversely:

PROPOSITION 2. Let \(f\) be a mapping of \(\mathbf{K}\) to \(\mathbf{R}_+\) belonging to \(\mathcal{V}(\mathbf{K})\) ; if there exists \(\mathbf{C} > 0\) such that \(f(n,1) \leqslant \mathbf{C}.n\) for every integer \(n > 0\) , \(f\) is an absolute value on \(\mathbf{K}\) .

By induction on r > 0 we deduce from \((\mathbf{U}_{\mathbf{A}})\) the relation

\[
f (x _ {1} + x _ {2} + \dots + x _ {2 ^ {r}}) \leqslant A ^ {r} \sup _ {1 \leqslant i \leqslant 2 ^ {r}} f (x _ {i})\tag{1}
\]

for every family \((x_{i})\) of \(2^{r}\) elements of \(\mathbf{K}\) . We set \(n = 2^{r} - 1\) ; for all \(x \in \mathbf{K}\) , we deduce from (1)

\[
(f (1 + x)) ^ {n} = f ((1 + x) ^ {n}) = f \left(\sum_ {i = 0} ^ {n} \binom {n} {i} x ^ {i}\right) \leqslant A ^ {r} \sup \left(f \left(\binom {n} {i}\right) (f (x)) ^ {i}\right)
\]

\[
\leqslant \mathrm{CA} ^ {r} \sum_ {i = 0} ^ {n} \binom {n} {i} (f (x)) ^ {i} = \mathrm{CA} ^ {r} (1 + f (x)) ^ {n}
\]

for \(f\left(\binom{n}{i}\right) \leqslant C\binom{n}{i}\) ; therefore

\[
f (1 + x) \leqslant \mathrm{C} ^ {1 / n} \mathrm{A} ^ {r / n} (1 + f (x)).
\]

Letting r tend to \(+\infty\) , we obtain \(f(1+x)\leqslant1+f(x)\) for all \(x\in K\) ; applying this inequality with x replaced by \(xy^{-1}\) (where \(y\neq0\) ) and taking account of \((\mathrm{VA}_{\mathrm{II}})\) , we obtain relation \((\mathrm{VA}_{\mathrm{III}})\) , which proves the proposition.

COROLLARY 1. For a mapping f from K to \(R_{+}\) to be an absolute value, it is necessary and sufficient that it satisfy conditions (VA \(_{I}\) ), (VA \(_{II}\) ) and (U \(_{2}\) ).

It is necessary, for \((\mathrm{VA}_{\mathrm{II}})\) implies

\[
f (x + y) \leqslant f (x) + f (y) \leqslant 2 \sup (f (x), f (y)).
\]

Conversely, suppose \(f\) satisfies \((\mathrm{VA}_{\mathrm{I}}), (\mathrm{VA}_{\mathrm{II}})\) and \((\mathrm{U}_2)\) ; for every integer \(n > 0\) , let \(r\) be the least integer such that \(2^r \geqslant n\) ; if in (1) A is replaced by 2, the \(x_i\) of index \(i \leqslant n\) by 1 and the \(x_i\) of index \(i > n\) by 0, we obtain

\[
f (n. 1) \leqslant 2 ^ {r} <   2 n;
\]

then Proposition 2 may be applied with \(\mathbf{C} = 2\) and hence \(f\) is an absolute value.

COROLLARY 2. For a mapping \(f\) of \(\mathbf{K}\) to \(\mathbf{R}_+\) to belong to \(\mathcal{V}(\mathbf{K})\) , it is necessary and sufficient that it be of the form \(g^t\) , where \(t > 0\) and \(g\) is an absolute value on \(\mathbf{K}\) .

To say that \(f\) satisfies \((\mathbf{U}_{\mathbf{A}})\) is equivalent to saying that \(f^s\) satisfies \((\mathbf{U}_{\mathbf{A}^s})\) ; as there exists \(s > 0\) such that \(\mathbf{A}^s \leqslant 2\) , Corollary 1 shows that for such a value of \(s, f^s\) is an absolute value.

